% Section 1.6, from lectures/L01/appendix/e_exercises.md.

\section{Exercises}
\label{c1:sec:exercises}

\begin{aside}
\textbf{How to check your work.} Exercises \ref{c1:ex:firstsub} and \ref{c1:ex:program} are checked by
this chapter's test suite: write the file at the path the specification gives, then run
\code{make test} from the repository root. Nothing has to be registered; creating the file is what
switches its tests on.

The rest will have worked solutions in the companion repository, published later as
\file{lectures/L01/appendix/f_solutions.md}, and the hand calculations have short answers in
Appendix~\ref{app:answers}. Work each one before you read it. Where an exercise
has a plausible wrong answer, the solution gives that too, and says why
it is wrong, because being told the right answer teaches you less than being shown the one you
nearly gave.

\textbf{No hardware is needed for any of this.} Everything is written, assembled, simulated and
measured on your laptop.
\end{aside}

Each exercise is labelled with its kind. There is exactly one \textbf{Cross-check} per chapter,
and in this one it is Exercise~\ref{c1:ex:crosscheck}.

% ------------------------------------------------------------------------------------------
\exercise{The upper half}{Recall}
\label{c1:ex:upper}
\begin{parts}
  \item State which registers \code{ldi} can name, and which it cannot.
  \item Explain \emph{why}, from the instruction's encoding rather than from the datasheet's
    table. Your answer should mention a specific number of bits.
  \item You are writing a routine that needs six working registers, and you have used \code{r16}
    to \code{r21}. A seventh value is a constant you need to load. Somebody suggests \code{r10},
    which is free. What does that cost you, in instructions and in cycles, every time you set it
    up?
  \item Name two instructions other than \code{ldi} that have the same restriction, and one
    instruction that does not.
\end{parts}

% ------------------------------------------------------------------------------------------
\exercise{Assembling four instructions by hand}{Hand calculation}
\label{c1:ex:assemble}
Using the encoding in \secref{c1:sec:encoding}, and nothing else, work out the 16-bit word each of
these assembles to. Show the four nibbles for at least the first one.
\begin{parts}
  \item \code{ldi r16, 0x2A}
  \item \code{ldi r24, 0x2A}
  \item \code{ldi r31, 0xFF}
  \item \code{ldi r20, 0x03}
  \item One of these is not a legal instruction: \code{ldi r15, 0x2A}. Say what the assembler does
    with it, and what would go wrong if it silently produced a word anyway.
\end{parts}

\noindent\textbf{Check yourself:} the encoding is \code{1110 KKKK dddd KKKK} with
\code{d = Rd - 16}, and the assembler will settle it. Put the instructions in a \code{.asm} file,
assemble it, and run \code{avr-objdump -D -m avr:5 -b ihex} over the hex; the bytes come out
reversed, for the reason in \secref{c1:sec:disasm}.

% ------------------------------------------------------------------------------------------
\exercise{Which space, which instruction}{Recall}
\label{c1:ex:spaces}
For each of the following, say which address space it names, what is stored there, and which
instruction or instructions can reach it.
\begin{parts}
  \item Data space \code{0x0005}
  \item Data space \code{0x0025}
  \item I/O address \code{0x05}
  \item Data space \code{0x0100}
  \item Program word address \code{0x0006}
  \item EEPROM address \code{0x0000}
  \item Two of the above name the same physical register. Which two, and what happens if you use
    one of the two numbers with an instruction expecting the other?
\end{parts}

% ------------------------------------------------------------------------------------------
\exercise{Choosing the branch}{Design}
\label{c1:ex:branch}
\code{r19} holds a loop counter and \code{r24} holds a limit. Both are unsigned.
\begin{parts}
  \item You want to leave the loop when the counter has reached the limit. Write the two
    instructions, and say which flag the second one reads.
  \item You want to stay in the loop while the counter is strictly below the limit. Write the two
    instructions.
  \item Somebody inserts an \code{inc r19} between the comparison and the branch. Explain what
    breaks, and why the symptom is a branch that goes the wrong way rather than an assembler
    error.
  \item Somebody inserts a \code{mov r20, r19} in the same place instead. Explain why this one is
    harmless, and say where you would look it up rather than guessing.
\end{parts}

% ------------------------------------------------------------------------------------------
\exercise{Your first subroutines}{Code}
\label{c1:ex:firstsub}
Write \code{shift_bits} and \code{shift_bits_inverted} in \file{drivers/source/utils.asm}, to the
specification in \secref{c1:sec:first}.
\begin{parts}
  \item Write both. Use \code{r18} and \code{r19} as scratch, for the reason in
    \secref{c1:sec:whyr18}.
  \item Assemble with \code{make build}, then disassemble your own routine with
    \code{avr-objdump -D -m avr:5 -b ihex drivers/build/drivers.hex} and read it against what you
    wrote. Both flags are required, for the reason in \secref{c1:sec:disasm}: a hex file says nothing
    about its own architecture. At least one line will have come back as a different instruction
    than you typed. Say which, and why that is not a mistake.
  \item Run \code{make test}.
  \item Before running anything else: predict what \code{shift_bits} returns for an argument of 8,
    and say what a reader thinking in C would predict instead.
\end{parts}

% ------------------------------------------------------------------------------------------
\exercise{A program that runs on its own}{Code}
\label{c1:ex:program}
Write \file{drivers/app/main.asm} to the specification in \secref{c1:sec:program}.
\begin{parts}
  \item Write the reset vector, the stack pointer initialisation, a call to \code{shift_bits}, the
    two instructions that light PB5, and the final loop.
  \item \code{make build} should now report that it built \file{app.hex} as well as
    \file{drivers.hex}, and \code{make test} should report two more passing tests than it did
    before.
  \item The hardware sets \code{SP} to \code{RAMEND} at reset by itself, so on a cold start your
    four instructions change nothing. Give the case where they do change something, and say what
    would go wrong without them.
  \item Remove the final loop and describe, without running it, what the machine does when it
    reaches the end of your code. What is in flash after your last instruction?
  \item Swap the two \code{sbi} instructions, so that \code{PORTB} is written before \code{DDRB}.
    The LED still ends up lit. Say what the pin is doing during the one instruction in between,
    and why that is not the same thing as being an output that happens to be low.
  \item You wrote \code{sbi DDRB, 5}, and the test that checks it reads data space address
    \code{0x24}. Both numbers are right. Explain, and give the one instruction that would write the
    same register using \code{0x24} directly.
\end{parts}

% ------------------------------------------------------------------------------------------
\exercise{Counting by hand}{Hand calculation}
\label{c1:ex:counting}
Using only the table in \secref{c1:sec:costs} and the rule in \secref{c1:sec:branchcost}, and without
running anything:
\begin{parts}
  \item Count the cycles \code{shift_bits} takes for an argument of 0. Show your working as a line
    per instruction, with how many times each one runs.
  \item Do the same for an argument of 7.
  \item From those two, write down the general formula in terms of $n$.
  \item How many times does the \code{cp} execute when $n$ is 3? How many times does the
    \code{breq}? How many of those \code{breq} executions are taken? These three numbers are not
    all the same, and getting them confused is the usual reason a hand count comes out wrong.
  \item Do the same for \code{shift_bits_inverted}, and state the difference between the two
    formulas in one sentence about a single instruction.
\end{parts}

% ------------------------------------------------------------------------------------------
\exercise[fillaccent]{What a subroutine costs}{Cross-check}
\label{c1:ex:crosscheck}
\emph{Compute it by hand, measure it, reconcile.}

This is the exercise the chapter exists for. Do the parts in order, and write each answer down
before starting the next; the point is lost if you look at the measurement first.
\begin{parts}
  \item \textbf{By hand.} From the exercise above, you have a formula. Evaluate it for $n = 0$ and
    for $n = 7$, and convert both to microseconds at 16\,MHz. One cycle is 62.5 nanoseconds, and
    doing that conversion in your head once is worth more than a function that does it for you.
  \item \textbf{The table, in full.} Write out the count for $n = 7$ in the shape
    \secref{c1:sec:onpaper} gives: one row per instruction, its cost, how many times it runs, and the
    product. You will need it in (d), and a formula you cannot decompose is a formula you cannot
    find the mistake in.
  \item \textbf{By measurement.}
\begin{shell}
make measure SYMBOL=shift_bits ARG=0
make measure SYMBOL=shift_bits ARG=7
\end{shell}
  \item \textbf{Reconcile the two.} They should agree exactly. If they do not, the table is more
    likely wrong than the simulator, and the row to check first is the \code{cp}, which runs
    $n + 1$ times rather than $n$. Find the mistake before continuing rather than adjusting the
    formula until it matches.
  \item \textbf{Now the call site.} Your \file{main.asm} does not just execute \code{shift_bits};
    it loads an argument and calls it. Count, by hand, the cycles of those two extra instructions,
    and write down what a loop that called \code{shift_bits(7)} over and over would cost \emph{per
    iteration}.
  \item \textbf{The discrepancy.} Your answer to (e) is larger than your answer to (a). State the
    difference in cycles, and state it again as a percentage of the (a) figure, once for $n = 0$
    and once for $n = 7$. \textbf{The two percentages are very different.} Explain why, in one
    sentence, without using the word ``overhead''.
  \item \textbf{What it means for a driver.} Chapter~2's \code{led_on} calls \code{shift_bits}
    with a pin number. Using your formula, how many times longer does \code{led_on} take for
    Arduino pin 13 than for Arduino pin 8? (Pin 13 is bit 5 of port B; pin 8 is bit 0 of port B.)
    Say whether that difference would matter, and give one situation in which it would.
  \item \code{make measure} reports a figure that does not include the \code{rcall}. Is that a
    defect in the tool? Give the case for and against, in a sentence each, and say what you would
    have it do instead if you disagree with the choice.
\end{parts}
